% Pulsar-M design decisions log -- included by pulsar-m.tex.
% Mirrors docs/design-decisions.md (now deleted; this file is the
% canonical LaTeX form). Each entry: decision, rationale,
% alternatives, status, date.

\subsection{DD-001 Target FIPS 204 output interchangeability (Class
N)}

\paragraph{Decision.}
Pulsar-M aims for Class N1 + N4 (NIST MPTC IR 8214C):
threshold-produced signatures verify under unmodified FIPS~204
ML-DSA.Verify.

\paragraph{Rationale.}
Output interchangeability gives the strongest NIST positioning.
FIPS-validated ML-DSA modules in 2026 (BoringSSL FIPS, AWS-LC,
OpenSSL 3.0) consume Pulsar-M certs without code changes. This is
the headline pitch for MPTC.

\paragraph{Alternative.}
Class S (special threshold-friendly primitive). Easier spec, but
does not give the FIPS-module-compatibility win. Pulsar (R-LWE)
takes this path.

\paragraph{Status.}
Targeted. Final class decision deferred until reference impl
demonstrates byte-equal output on at least one parameter set.

\subsection{DD-002 Hash family: SHA-3 (cSHAKE256 / KMAC256 /
TupleHash256) only}

\paragraph{Decision.}
Pulsar-M's NIST profile uses exclusively the FIPS~202 + SP 800-185
hash family. No BLAKE3.

\paragraph{Rationale.}
NIST MPTC \S4.6 lists allowed symmetric primitives for Class N: AES,
Ascon, SHA-2, SHA-3, SHAKE/cSHAKE, HMAC/KMAC/CMAC/GMAC. BLAKE3 is
not on the list. Including BLAKE3 in the NIST submission profile is
a strategic liability.

The same SHAKE-256 sponge is used internally by ML-DSA-65 (FIPS~204
\S5.1). Sharing the hash family with the underlying primitive
removes a whole class of cryptanalytic concern (cross-domain
collisions, mixed hash-family soundness arguments).

\paragraph{Alternative considered.}
Dual-suite with BLAKE3 as the production profile and SHA-3 as the
NIST profile, like Pulsar's current dual-suite. Rejected because
Pulsar-M is NIST-track and shouldn't ship a non-NIST default.

\paragraph{Status.}
Final.

\subsection{DD-003 Lattice: Module-LWE per FIPS 204, not Ring-LWE}

\paragraph{Decision.}
Pulsar-M operates over $\Rq^k$ (module of polynomials) with the same
$(q, k, \ell, \eta, \beta, \omega, \ldots)$ parameter set as
ML-DSA-65. Not $\Rq$ (single polynomial) as Pulsar / Ringtail use.

\paragraph{Rationale.}
This is the whole point of the project. Class N1 interchangeability
requires the protocol's signing operation to land on exactly the
FIPS~204 $\sigma$ encoding, which is M-LWE-shaped.

\paragraph{Alternative.}
Stay R-LWE. Easier (we already have Pulsar). But then we lose the
Class N positioning and become Class S like Pulsar.

\paragraph{Status.}
Final.

\subsection{DD-004 Reference implementation language: Go}

\paragraph{Decision.}
Reference implementation is Go (\texttt{ref/go/}). C implementation
comes after spec encodings freeze (\texttt{ref/c/}).

\paragraph{Rationale.}
Go matches Pulsar's reference; reuses our Pedersen DKG and reshare
patterns from \texttt{pulsar/dkg2/} and \texttt{pulsar/reshare/}.
Clear, boring, no assembly. NIST MPTC \S4.3 doesn't mandate
language.

\paragraph{Alternative.}
C reference first. Standard NIST PQC submission style. Rejected for
round 1 because we need to ship by 2026-Nov-16; Go is faster to
write and review. C reference becomes a follow-on artifact for round
2 review.

\paragraph{Status.}
Final for MPTC round 1.

\subsection{DD-005 DKG: Pedersen-style over $\Rq^k$}

\paragraph{Decision.}
Distributed keygen uses Pedersen DKG over the M-LWE module ring,
with Feldman-style verifiable secret sharing replaced by the
module-shaped Pedersen commitment scheme.

\paragraph{Rationale.}
Pulsar's \texttt{dkg2/} package already does this for $\Rq$ (R-LWE).
The M-LWE port preserves the soundness argument with constant
overhead per coefficient slot.

\paragraph{Alternative.}
Trusted-dealer DKG (academic Ringtail). Rejected: incompatible with
open public chains (HIP-0077 \S``Pulsar vs Ringtail'').

\paragraph{Status.}
Final.

\subsection{DD-006 Proactive resharing: HJKY97-style with
beacon-randomized quorum}

\paragraph{Decision.}
Long-lived deployments use proactive secret resharing with a
beacon-derived per-reshare quorum selection (not deterministic
smallest-ID).

\paragraph{Rationale.}
Deterministic quorum selection (Pulsar's current
\texttt{reshare.go:206-209}) lets a mobile adversary game the
corrupt-set positioning. HIP-0077 red review F10 flagged this.
Pulsar-M's reshare selects via the chain's randomness beacon at the
reshare epoch boundary.

\paragraph{Alternative.}
Deterministic quorum, accept the leakage as a documented limitation.
Rejected for MPTC submission.

\paragraph{Status.}
Final.

\subsection{DD-007 No application-level logging in secret-touching
code paths}

\paragraph{Decision.}
Secret-touching functions (Sign, Verify-of-shares, DKG verification,
reshare verification) emit zero log lines. Constant-time analysis
treats any log call as a failed gate.

\paragraph{Rationale.}
Pulsar's pre-fix \texttt{log.Println(Bsquare, sumSquares)} in
\texttt{sign.go} (HIP-0077 red F8) is exactly the failure mode this
rule prevents. A logged value derived from secret state is a side
channel even at debug level.

\paragraph{Alternative.}
Allow debug logging gated by build tag. Rejected: build tags get
accidentally enabled in production; the rule is absolute.

\paragraph{Status.}
Final. CI gate: \texttt{go vet -vettool=...} plus a custom linter
that rejects \texttt{log.*}, \texttt{fmt.Print*}, \texttt{panic} in
\texttt{pkg/sign/}, \texttt{pkg/dkg/}, \texttt{pkg/reshare/}.

\subsection{DD-008 Encoding freeze before C reference}

\paragraph{Decision.}
Reference Go locks every byte of the wire encoding before the C
reference begins. Inputs to the freeze: spec \S``Encodings'', KAT
suite v0, cross-vendor review with at least one external lattice
impl team.

\paragraph{Rationale.}
Two implementations diverging on encoding is a recurring source of
NIST PQC submission rejections. Freeze first, port second.

\paragraph{Status.}
Targeted for end of August 2026.

\subsection{DD-009 Patent claims: prepared in advance}

\paragraph{Decision.}
Patent-claim disclosures collected in
\texttt{spec/patent-notes.tex} from day 1 of the repo. Every
contribution adds the contributor's patent-claim notice (or explicit
``no claims'').

\paragraph{Rationale.}
NIST MPTC \S6 requires Notes on Patent Claims as a package
deliverable. Late collection blows up the submission timeline.

\paragraph{Status.}
In progress.

\subsection{DD-010 Two-track NIST plan: MPTC now, CMVP later}

\paragraph{Decision.}
Track A targets MPTC submission 2026-Nov-16. Track B prepares for
eventual NIST PQC standardization + CMVP module validation, which is
a multi-year process not scheduled.

\paragraph{Rationale.}
MPTC is the active NIST process for threshold schemes right now.
CMVP requires an approved standard, which threshold schemes don't
yet have. We can prepare a clean module boundary for future CMVP
work without waiting for the standard.

\paragraph{Status.}
Track A active. Track B is a parallel readiness exercise, not a
near-term gate.
